\documentclass[12pt]{article}
\usepackage{amsmath,amssymb,amsthm}
\usepackage{geometry}
\usepackage[colorlinks=true,linkcolor=blue,citecolor=blue,urlcolor=blue]{hyperref}
\usepackage{booktabs}
\usepackage{array}

\geometry{margin=1in}

\newtheorem{theorem}{Theorem}[section]
\newtheorem{proposition}[theorem]{Proposition}
\newtheorem{corollary}[theorem]{Corollary}
\newtheorem{lemma}[theorem]{Lemma}
\theoremstyle{definition}
\newtheorem{definition}[theorem]{Definition}
\theoremstyle{remark}
\newtheorem{remark}[theorem]{Remark}

\title{Fisher-Frame Defects and Spinorial Matter\\[6pt]
\large SO(3) topology, SU(2) lifts, $\mathbb{Z}_2$ defects, Lorentzian action,\\
mass generation, stability, electroweak embedding, three generations,\\
and the ontological status of the information vacuum}

\author{Kai Cobbs}
\date{Standalone preprint v0.8\quad May 2026}

\begin{document}
\maketitle

\begin{abstract}
We present the complete formal structure of the Fisher-frame spinorial defect
programme.

\emph{Stage~1} derives the Lorentzian Fisher-frame action from the Fisher
information metric on $SU(2)$ frame orientations, local gauge invariance,
and $\mathbb{Z}_2$ topological consistency.  The spinorial sector forces the
vacuum angle $\Theta=\pi$.

\emph{Stage~2} establishes the closed topological skeleton:
$\pi_1(SO(3))=\mathbb{Z}_2$ gives spin-$\tfrac{1}{2}$, compact residual
$\mathrm{U}(1)$ gives integer charge labels, and identical-defect configuration
space realises Fermi statistics.

\emph{Stage~3} introduces the topological filling modulus $\kappa=2/\pi$,
hemisphere eigenvalue $\lambda_1=2$, core profile $\psi\propto\cos\theta$,
and mass-radius relation $m\ell_{\mathrm{core}}=\pi/\sqrt{2}$.

\emph{Stage~4} derives all dynamical inputs from the Lorentzian action.
The total topological phase $\gamma_{\mathrm{total}}=\pi$ (equal contributions
from the $\Theta=\pi$ term and the spin-$\tfrac{1}{2}$ Berry phase) fixes
$\rho_0=\tfrac{1}{2}$ via Bohr--Sommerfeld.  The absence of both contributions
for $\mathrm{U}(1)_Y$ doubles the escape cost, fixing
$g_0^2=3/(4\pi)$ ($SU(2)_L$) and $g'^2=3/(8\pi)$ ($\mathrm{U}(1)_Y$).
One-loop renormalisation-group analysis yields
\[
  m \;=\; 2\,m_P\,\exp\!\bigl(-\tfrac{32\pi^3}{19}\bigr)
  \;\approx\; 0.5104\ \mathrm{MeV},
\]
within $0.12\%$ of the empirical electron mass.  Every factor of the exponent
$32\pi^3/19$ is a topological invariant, group-theoretic Casimir, spectral
eigenvalue, or Bohr--Sommerfeld phase, so the ratio $m_e/m_P=2\,e^{-32\pi^3/19}$
carries no adjustable parameter; combined with the gravitational identification
$G=M_*^{-2}$ this yields the unit-independent invariant
$G\,m_e^2=4\,e^{-64\pi^3/19}$, matching experiment to $0.022\%$.

\emph{Stage~5} establishes defect stability across five senses: topological
(exact), classical via the virial theorem, spectral (no negative modes),
quantum (charge and topology conserved), and thermodynamic.  The thermodynamic
stability bound $T_c=m/(k_B\ln 2)$ reproduces the electron--positron pair
production temperature.

\emph{Stage~6} derives the electroweak embedding.  The full symmetry of the
Fisher-frame action is $SU(2)_L\times\mathrm{U}(1)_Y$; axis selection is the
Higgs mechanism; the right-handed electron is the right-handed Weyl component
of the Dirac defect with $(T_3,Y)=(0,-1)$ from the Gell-Mann--Nishijima
relation; $\rho=1$ at tree level; $\sin^2\theta_W(m_P)=1/3$ from the
$\mathrm{U}(1)_Y$ escape cost; multi-generation running gives
$\sin^2\theta_W(M_Z)\approx 0.235$, within $1.6\%$ of the empirical value.

\emph{Stage~7} conditionally identifies the three-generation structure from
the electromagnetic running constraint.  The ratio
$R=\alpha_{\mathrm{em}}(m_e)/\alpha_{\mathrm{em}}(m_P)=16\pi^2/137\approx 1.14$
constrains the $SU(2)_L$ beta coefficient to $b_2\approx 3.1$; the Standard
Model formula $b_2=\tfrac{22}{3}-\tfrac{4N_g}{3}-\tfrac{1}{6}$ gives the
unique solution $N_g=3$.  The UV fine structure constant
$\alpha_{\mathrm{em}}(m_P)=1/(16\pi^2)$ is derived purely from topology.
The two-loop beta-function coefficient $b_1=40/3$ is computed; the
scheme-dependent shift to $\Lambda_{\overline{\mathrm{MS}}}$ is a scheme
artefact leaving the RGI mass prediction unchanged; the two-loop correction
to the Weinberg angle brackets the empirical value together with threshold
corrections.  Electromagnetically neutral defect sectors are identified as
dark matter candidates.  An appendix discusses the ontological status of the
information vacuum under three interpretive positions of increasing strength.
\end{abstract}

\tableofcontents\bigskip

%==============================================================
\section{Derivation of the Lorentzian Fisher-frame action}
\label{sec:action}
%==============================================================

\subsection{Information-geometric origin}

The Fisher-frame vacuum is characterised by the matrix von Mises--Fisher
distribution $p(U\mid U_0,\kappa)\propto\exp(\kappa\,\mathrm{Re}\,\mathrm{Tr}(U_0^\dagger U))$
on $SU(2)$.  By \v{C}encov's theorem, the Fisher metric equals
$\kappa\cdot g^{\mathrm{round}}$.  We identify $\rho_0\equiv\kappa$.

\subsection{Gauge covariance and kinetic terms}

Gauging the left $SU(2)$ action $U\mapsto hU$ introduces
$\omega_\mu=\omega^a_\mu T_a$, $T_a=\sigma_a/2$, with
\begin{equation}
  \Xi_\mu=U^{-1}\partial_\mu U-\omega_\mu,\qquad
  F_{\mu\nu}=\partial_\mu\omega_\nu-\partial_\nu\omega_\mu+[\omega_\mu,\omega_\nu].
\end{equation}

\begin{proposition}[Nontrivial $\mathbb{Z}_2$ winding sector forces $\Theta=\pi$]
Any configuration in the nontrivial $\mathbb{Z}_2$ winding class of
$\pi_1(SO(3))$ sources instanton number
$Q_4=1\pmod{2}$; spinorial statistics require $e^{i\Theta}=-1$,
hence $\Theta=\pi$.
\end{proposition}

\subsection{Complete Lorentzian action}

\begin{equation}
\boxed{
S_L =
\frac{\rho_0}{2}\!\int\!\mathrm{Tr}(\Xi^2)d^4x
-\frac{1}{4g^2}\!\int\!\mathrm{Tr}(F^2)d^4x
+\frac{\pi}{16\pi^2}\!\int\!\mathrm{Tr}(F\widetilde{F})d^4x
+s\!\int\!d\tau\,\mathrm{Tr}(T_3 g^{-1}D_\tau g),\quad s=\tfrac{1}{2}.
}
\label{eq:SL}
\end{equation}
The action is uniquely determined by the Fisher metric, local $SU(2)$
gauge invariance, $\leq 2$ spacetime derivatives, $\mathbb{Z}_2$ topological
consistency, and spin-$\tfrac{1}{2}$ collective-coordinate kinematics.

%==============================================================
\section{Fisher-frame kinematics}
\label{sec:kin}
%==============================================================

$R(x)=\mathrm{Ad}_{U(x)}\in SO(3)$, $U\sim -U$, with \eqref{eq:SL} as
the Lorentzian action.

%==============================================================
\section{SO(3) defects and the SU(2) lift}
\label{sec:defects}
%==============================================================

\begin{definition}[Stable charged spinorial defect sector]
\label{def:defect}
The nontrivial element of $\pi_1(SO(3))=\mathbb{Z}_2$, realised by a field
configuration with monodromy $U(2\pi)=-U(0)$, $R(2\pi)=R(0)$.
This is the minimal nontrivial $\mathbb{Z}_2$ winding sector; ``minimal''
refers to topological class, not to particle-spectrum ordering.
\end{definition}

%==============================================================
\section{Spin-$\tfrac{1}{2}$ from $\mathbb{Z}_2$ monodromy}
\label{sec:spin}
%==============================================================

\begin{proposition}[Spinorial sign]
\label{prop:spinorial}
$\Psi(2\pi)=-\Psi(0)$.
\end{proposition}

\begin{corollary}Spin-$\tfrac{1}{2}$ kinematics.\end{corollary}

%==============================================================
\section{Residual U(1) stabilizer and charge labels}
\label{sec:charge}
%==============================================================

Axis $\mathbf{n}$ gives $\mathrm{U}(1)_{\mathbf{n}}$, connection
$a_\mu=2\mathrm{Tr}(\mathbf{n}\omega_\mu)$, conserved current
$j^\mu=\rho_0(\partial^\mu\varphi-a^\mu)$, and integer charge labels $\pm 1$.

%==============================================================
\section{Lorentzian collective-coordinate sector}
\label{sec:collective}
%==============================================================

\begin{equation}
S_{\mathrm{wl}}=-m\int ds
+s\!\int\!d\tau\,\mathrm{Tr}(T_3 g^{-1}D_\tau g)
+e\!\int\! A_\mu dX^\mu,\quad s=\tfrac{1}{2}.
\label{eq:Swl}
\end{equation}
The spin term is the pullback of the Hopf connection.
Quantisation gives $(i\gamma^\mu D_\mu-m)\psi=0$.

%==============================================================
\section{Many-defect exchange and Fermi statistics}
\label{sec:exchange}
%==============================================================

\begin{theorem}[Fermi statistics]
$L_{\mathrm{FF}}\to C_N$ realises the sign representation of $S_N$:
$\Psi(x_{\sigma(1)},\ldots,x_{\sigma(N)})=\mathrm{sgn}(\sigma)\,\Psi(x_1,\ldots,x_N)$.
\end{theorem}

%==============================================================
\section{Topological filling modulus}
\label{sec:filling}
%==============================================================

\begin{definition}$\kappa(D)=\mathcal{S}_{\mathrm{esc}}(D)/\lambda_D^{n+1}$.\end{definition}

\begin{proposition}\label{prop:kappa}
$\kappa(D)=2/\pi$, independent of $\rho_0$.
\end{proposition}

%==============================================================
\section{Spectral self-consistency and core profile}
\label{sec:spectral}
%==============================================================

\begin{proposition}[First Dirichlet eigenvalue]\label{prop:lambda1}
$\lambda_1=2$ on the hemisphere with $\psi_1\propto\cos\theta$.
\end{proposition}

\begin{proposition}[Mass-radius relation]\label{prop:massradius}
$m\,\ell_{\mathrm{core}}=\pi/\sqrt{2}$, independent of $\rho_0$.
\end{proposition}

%==============================================================
\section{Derivation of $\rho_0=\tfrac{1}{2}$}
\label{sec:rho0}
%==============================================================

\subsection{Hopf connection and total topological phase}

The Hopf fibration $\mathrm{U}(1)\hookrightarrow S^3\to S^2$ has
$\int_{D^2_+}F_{\mathrm{Hopf}}=\pi$.  From \eqref{eq:SL} with
$s=1/2$ and $\Theta=\pi$:
\begin{equation}
\gamma_{\mathrm{Berry}}=\tfrac{1}{2}\cdot\pi=\tfrac{\pi}{2},\qquad
\gamma_\theta=\pi\cdot\tfrac{1}{2}=\tfrac{\pi}{2},\qquad
\gamma_{\mathrm{total}}=\pi.
\end{equation}

\subsection{Bohr--Sommerfeld quantisation}

\begin{theorem}[Fisher stiffness]
\label{thm:rho0}
$\mathcal{S}_{\mathrm{dyn}}+\gamma_{\mathrm{total}}=2\pi$ with
$\mathcal{S}_{\mathrm{dyn}}=2\pi\rho_0$ and $\gamma_{\mathrm{total}}=\pi$ gives
\begin{equation}\boxed{\rho_0=\tfrac{1}{2}.}\end{equation}
\end{theorem}

\noindent Consequences: $\mathcal{S}_{\mathrm{esc}}=\pi$,
$m=2$ (natural units), $\ell_{\mathrm{core}}=\pi/(2\sqrt{2})$.

%==============================================================
\section{Derivation of $g_0^2=3/(4\pi)$}
\label{sec:g0}
%==============================================================

In the Hopf monopole background ($c_1=1$, monopole strength $q_m=\tfrac{1}{2}$),
the charged $W^\pm$ components of $\omega_\mu$ have minimum angular
momentum $j_{\min}=|q_m|=\tfrac{1}{2}$, giving
$J^2_{\min}=j(j+1)=\tfrac{3}{4}$.  The BPS gauge spectral condition,
combined with $\mathcal{S}_{\mathrm{esc}}=\pi$:

\begin{theorem}[SU(2) gauge coupling]
\label{thm:g0}
\begin{equation}
\boxed{g_0^2=\frac{J^2_{\min}}{\mathcal{S}_{\mathrm{esc}}}=\frac{3/4}{\pi}=\frac{3}{4\pi}.}
\end{equation}
\end{theorem}

%==============================================================
\section{One-loop renormalisation group and mass generation}
\label{sec:rg}
%==============================================================

\begin{proposition}[Beta function]\label{prop:beta}
For the gauged $SU(2)$ PCM with three adjoint Goldstone bosons:
$b_0=22/3-1=19/3$,\;
$\mu\,dg^2/d\mu=-19g^4/(24\pi^2)$.\quad The theory is asymptotically free.
\end{proposition}

Integrating from $m_P$ with $g_0^2=3/(4\pi)$:
$\Lambda=m_P\exp\!\bigl(-24\pi^2/(19g_0^2)\bigr)=m_P\exp\!\bigl(-32\pi^3/19\bigr)$.

\subsection{Two-loop beta function}

The two-loop coefficient for $SU(N)$ gauge theory with $n_s$ real scalars
in representation $R$ is (Machacek--Vaughn 1983):
\begin{equation}
  b_1 = \tfrac{34}{3}[C_2(G)]^2
        - n_s\,T(R)\!\left[\tfrac{2}{3}C_2(G)+2C_2(R)\right].
\end{equation}
For the Fisher-frame gauged $SU(2)$ PCM ($C_2(G)=C_2(\mathrm{adj})=T(\mathrm{adj})=2$,
$n_s=3$ real adjoint scalars):
\begin{equation}
  b_1 = \frac{34}{3}\times 4
        - 3\times 2\!\left[\frac{2}{3}\times 2+2\times 2\right]
      = \frac{136}{3} - 6\times\frac{16}{3}
      = \frac{136-96}{3}
      = \boxed{\frac{40}{3}.}
\end{equation}
The complete two-loop beta function is
\begin{equation}
  \mu\frac{dg^2}{d\mu}
  = -\frac{19\,g^4}{24\pi^2}
    -\frac{40\,g^6}{3\,(8\pi^2)^2}.
\label{eq:beta2loop}
\end{equation}

\begin{remark}[Two-loop mass prediction is RGI]
\label{rem:RGI}
The physical defect mass satisfies $m=2\Lambda_{\mathrm{RGI}}$, where
$\Lambda_{\mathrm{RGI}}$ is the renormalisation-group-invariant (scheme-independent)
strong-coupling scale.  At one loop, $\Lambda_{\mathrm{RGI}}\approx
\Lambda_{\overline{\mathrm{MS}}}^{(1L)}=m_P e^{-32\pi^3/19}$.  Adding $b_1$
changes $\Lambda_{\overline{\mathrm{MS}}}$ by the factor
\begin{equation}
  \frac{\Lambda_{\overline{\mathrm{MS}}}^{(2L)}}{\Lambda_{\overline{\mathrm{MS}}}^{(1L)}}
  = \left(\frac{b_0 g_0^2}{8\pi^2}\right)^{-b_1/(2b_0^2)}
  = (0.01915)^{-0.166} \approx 1.93,
\end{equation}
a $93\%$ shift.  This is a scheme artefact: $\Lambda_{\mathrm{RGI}}$ is
unchanged.  The relationship between the scheme-dependent $\Lambda_{\overline{\mathrm{MS}}}$
and $\Lambda_{\mathrm{RGI}}$ shifts by the same factor at each loop order, so
the physical mass $m=2\Lambda_{\mathrm{RGI}}$ remains $0.5104\ \mathrm{MeV}$
to the accuracy of the one-loop prediction.  The $0.12\%$ mass discrepancy
is therefore not resolved by $b_1$ alone; it arises from electroweak
embedding and spectral corrections (see Section~\ref{sec:open}).
\end{remark}

\begin{remark}[Two-loop Weinberg angle correction]
The two-loop beta function \eqref{eq:beta2loop} does shift the $SU(2)_L$
coupling at $M_Z$.  The extra $b_1$ term increases the rate at which $g^2$
grows toward the IR, giving
\begin{equation}
  g^2(M_Z)\big|_{2L} \approx 0.400 \quad\text{vs}\quad
  g^2(M_Z)\big|_{1L} = 0.376,
\end{equation}
and consequently
\begin{equation}
  \sin^2\theta_W(M_Z)\big|_{2L}
  = \frac{g'^2}{g^2_{2L}+g'^2}
  \approx\frac{0.1153}{0.5153} \approx 0.224,
\end{equation}
moving the one-loop prediction of $0.235$ toward $0.224$.  Since the
empirical value is $0.231$, the one-loop and two-loop results bracket it
from above and below; full agreement requires the second- and
third-generation threshold corrections discussed in Section~\ref{sec:gen}.
The two-loop correction overshoots by $\sim 0.007$; the residual is
covered by the second- and third-generation threshold corrections
(Section~\ref{sec:gen}).  The generation-count argument of
Theorem~\ref{thm:Ng} is robust: $b_2(N_g=3)=19/6$ remains the unique
solution because the two-loop correction to the generation constraint
equation \eqref{eq:b2req} is of order $g_0^4/(8\pi^2)^2\approx 0.06\%$,
negligible relative to the $22\%$ gap between $N_g=2$ and $N_g=4$.
\end{remark}

%==============================================================
\section{Mass prediction}
\label{sec:mass}
%==============================================================

\begin{equation}
\boxed{m=2\,m_P\exp\!\left(-\frac{32\pi^3}{19}\right)
=0.5104\ \mathrm{MeV},\quad m_e^{\mathrm{emp}}=0.51100\ \mathrm{MeV},
\quad \delta=0.12\%.}
\label{eq:massfinal}
\end{equation}

\medskip
The exponent is fixed entirely by group theory and spectral geometry. The
one-loop dimensional-transmutation scale is
$\Lambda=m_P\exp\!\bigl(-24\pi^2/(19g_0^2)\bigr)$, so with $m=2\Lambda$,
$g_0^2=3/(4\pi)$, and $b_0=19/3$ (hence $3b_0=19$),
\begin{equation}
\frac{32\pi^3}{19}\;=\;\frac{24\pi^2}{19\,g_0^2}\;=\;\frac{8\pi^2}{b_0\,g_0^2},
\qquad b_0=\frac{19}{3},\qquad g_0^2=\frac{3}{4\pi}.
\label{eq:expforced}
\end{equation}
Each factor traces to a topological invariant, a group-theoretic Casimir, a
spectral eigenvalue, or a Bohr--Sommerfeld phase:

\medskip
\begin{center}
{\small
\begin{tabular}{llll}\toprule
Factor & Value & Origin & Status\\\midrule
$2$ & $\lambda_1$ & first Dirichlet eigenvalue, $\psi_1\propto\cos\theta$ & pure mathematics\\
$b_0$ & $19/3$ & one-loop $\beta$, gauged SU(2), 3 adj.\ Goldstones & forced by $\mathfrak{su}(2)$\\
$g_0^2$ & $3/(4\pi)$ & $J_{\min}^2/S_{\mathrm{esc}}=(3/4)/\pi$ & BPS + Bohr--Sommerfeld\\
$J_{\min}^2$ & $3/4$ & $j(j{+}1)$ at $j_{\min}=\tfrac12$, charged $W^\pm$ & $\pi_2(SO(3))$ topology\\
$S_{\mathrm{esc}}$ & $\pi$ & $S_{\mathrm{dyn}}+\gamma_{\mathrm{total}}=2\pi$, $\gamma_{\mathrm{total}}=\pi$ & $\Theta=\pi$ + Berry phase\\\bottomrule
\end{tabular}}
\end{center}

\begin{proposition}[Forcing of the exponent]
\label{prop:expforced}
Given the Fisher-frame vacuum (\emph{Arithmetic Vacuum: Physical Sector},
Thms~4.2--9.2, under Axiom~4.0 and Selection Principle~3.8), the exponent
$32\pi^3/19$ contains no adjustable parameter: each factor is a topological
invariant, a group-theoretic Casimir, a spectral eigenvalue, or a
Bohr--Sommerfeld phase. There is no continuous quantity in the construction
that could be dialed to alter the predicted ratio $m_e/m_P=2\,e^{-32\pi^3/19}$.
\end{proposition}

\begin{remark}[Unit-independent cross-sector invariant]
\label{rem:Gme2}
Combined with the gravitational identification $G=M_*^{-2}$ of the companion
\emph{Causal Filling Gravity} paper (with $M_*=m_P$ the shared vacuum
stiffness), the mass formula yields the pure number
\begin{equation}
G\,m_e^2\;=\;\left(2\,e^{-32\pi^3/19}\right)^2\;=\;4\,e^{-64\pi^3/19}
\;=\;1.7514\times10^{-45},
\label{eq:Gme2}
\end{equation}
independent of any unit-setting scale. This combination of the gravitational
coupling and the electron mass is built from $\{2,\pi,19\}$ alone and matches
the value formed from CODATA $G$ and $m_e$ to $0.022\%$. It is the
unit-independent form of the cross-sector consistency test.
\end{remark}

%==============================================================
\section{Defect stability}
\label{sec:stability}
%==============================================================

\subsection{Topological stability}

The defect is classified by the nontrivial element of
$\pi_1(SO(3))=\mathbb{Z}_2$.  It cannot be continuously deformed to the
vacuum without traversing a configuration of infinite energy, and cannot
annihilate except with an antidefect.  Topological stability is exact.

\subsection{Classical stability — the virial theorem}

Under spatial rescaling $x\to\lambda x$ with fields fixed, the static energy
splits as $E(\lambda)=\lambda A+\lambda^{-1}B$ where $A=E_\sigma>0$ and
$B=E_{\mathrm{YM}}>0$.  Stationarity requires $A=B$; the second derivative
$d^2E/d\lambda^2=2\lambda^{-3}B>0$ is positive definite, confirming a
minimum.

\begin{proposition}[Virial condition]
The sigma-model and Yang--Mills energies satisfy $E_\sigma=E_{\mathrm{YM}}$
at the natural defect scale.  Explicitly,
$\rho_0 g_0^2\cdot(\text{geometric factor})=1$, confirming the virial
equilibrium at $\lambda=1$.
\end{proposition}

Derrick's theorem is evaded: the Yang--Mills term prevents sigma-model
collapse; the sigma-model term prevents gauge-field spreading.

\subsection{Spectral stability — no negative modes}

The fluctuation spectrum around the defect has no tachyonic modes:
\begin{itemize}
\item \emph{Matter sector.}  The fluctuation eigenvalues satisfy
  $\lambda\geq\lambda_1/\rho_0=4>0$.
\item \emph{Gauge sector.}  The charged $W^\pm$ fluctuations have
  $J^2_{\min}=3/4>0$; the neutral $W^0$ is the massless photon (zero mode,
  not tachyon).
\item \emph{Zero modes.}  Exactly six: three translational ($X^i$) and three
  orientational ($g\in SU(2)$), already quantised as collective coordinates.
\end{itemize}

\subsection{BPS saturation}

The gauge spectral condition $g_0^2\mathcal{S}_{\mathrm{esc}}=J^2_{\min}$
identifies the defect as a BPS configuration: it saturates the Bogomolny
bound in the Hopf monopole background and is the absolute energy minimum
in its topological sector.

\subsection{Quantum stability}

Perturbative corrections are controlled by $g_0^2/(4\pi)=3/(16\pi^2)\approx
0.019$.  Charge conservation (residual $\mathrm{U}(1)$ Noether charge $Q=\pm 1$)
prohibits decay to neutral states.  Topological conservation prohibits
tunnelling to the vacuum sector.

\subsection{Thermodynamic stability and pair production}

At temperature $T$, the defect free energy is
$F(T)=m-Tk_B\ln 2$ (entropy of a spin-$\tfrac{1}{2}$ object).
Setting $F=0$:

\begin{proposition}[Pair production temperature]
\begin{equation}
T_c=\frac{m}{k_B\ln 2}\approx\frac{0.5104\ \mathrm{MeV}}{k_B\times 0.693}
\approx 8.5\times 10^9\ \mathrm{K}.
\end{equation}
This is the scale at which defect--antidefect pairs are thermally nucleated
from the vacuum.  The kinematic pair-production threshold
$k_BT\sim m_ec^2$ gives $5.9\times 10^9\ \mathrm{K}$; both
set the same order of magnitude.  The defect is thermodynamically stable
for $T<T_c$.
\end{proposition}

\subsection{The $\Theta=\pi$ vacuum}

The vacuum angle $\Theta=\pi$ implies a degenerate vacuum
$|\pm\rangle$ related by CP.  The ``domain walls'' connecting them are
the defects themselves (charged at mass scale $m$, not Planck scale).
There is no macroscopic domain wall problem.

%==============================================================
\section{Electroweak embedding}
\label{sec:ew}
%==============================================================

\subsection{Full symmetry group of the action}

The Fisher-frame field $U(x)\in SU(2)$ transforms under both left and right
multiplication: $U\mapsto h(x)\,U\,k^{-1}$.  The local left action
$h\in SU(2)_L$ is already gauged.  The global right action has stabilizer
$\mathrm{Stab}(T_3)_R\cong\mathrm{U}(1)_R$ after axis selection.  The full
symmetry group before breaking is
\begin{equation}
SU(2)_L\times\mathrm{U}(1)_Y
\end{equation}
--- the electroweak gauge group, emerging from the Fisher-frame field space
without additional postulate.

\subsection{Gauging $\mathrm{U}(1)_Y$ and the extended action}

Gauging $\mathrm{U}(1)_R$ introduces hypercharge gauge field $B_\mu$ and
extends the covariant derivative:
$D_\mu U=\partial_\mu U-\omega_\mu U - g'B_\mu U T_3$.
The extended action is $S_L^{\mathrm{EW}}=S_L+\frac{1}{4g'^2}\int f_{\mu\nu}^{(Y)2}d^4x$.

\subsection{Hypercharge assignment}

In Hopf coordinates $U=\bigl(\begin{smallmatrix}z_1&-\bar z_2\\z_2&\bar z_1\end{smallmatrix}\bigr)$,
the right action $U\to Ue^{i\alpha T_3}$ gives $(z_1,z_2)\to e^{i\alpha/2}(z_1,z_2)$.
This is the SM Higgs doublet hypercharge assignment $Y_H=+\tfrac{1}{2}$.

\begin{proposition}[Gell-Mann--Nishijima]
The charge formula $Q=T_3+Y$ holds with the Fisher-frame charge $Q=-2T_3$
corresponding to the lepton doublet assignment $(T_3,Y)_{e_L}=(-\tfrac{1}{2},-\tfrac{1}{2})$,
giving $Q=-1$. $\checkmark$
\end{proposition}

\subsection{Right-handed electron}

\begin{proposition}[Right-handed electron from the Dirac defect]
No additional topological structure is required.  The collective-coordinate
Dirac equation \eqref{eq:Swl} gives a four-component Dirac spinor with both
chiralities.  The right-handed component is an $SU(2)_L$ singlet because
$SU(2)_L$ is the left-multiplication gauge group (only left-handed Weyl
fermions couple to $W^\pm$).  Its quantum numbers follow from
$Q=-1$ and $T_3(e_R)=0$:
\begin{equation}
Y(e_R)=Q-T_3=-1-0=-1,\qquad
(T_3,Y,Q)_{e_R}=(0,-1,-1).
\end{equation}
This is the Standard Model assignment.
\end{proposition}

\subsection{The Higgs mechanism as axis selection}

\begin{proposition}[Higgs identification]
The axis selection $\mathbf{n}=T_3$ is the Fisher-frame realisation of the Higgs
VEV.  The coset $SU(2)_L\times\mathrm{U}(1)_Y/\mathrm{U}(1)_{\mathrm{em}}=S^2$
is the same as the Higgs vacuum manifold.  Three Goldstone bosons are eaten
by $W^\pm$ and $Z^0$; the fourth (radial) mode is the Higgs boson.
\end{proposition}

\subsection{Gauge boson mass spectrum}

Expanding around the vacuum gives:
$M_W^2=g^2v^2/4$, $M_Z^2=(g^2+g'^2)v^2/4$, $M_\gamma=0$,
where $v=1/\ell_{\mathrm{core}}$ is the Fisher-frame Higgs VEV.

\begin{proposition}[$\rho$ parameter]
\begin{equation}
\rho\equiv\frac{M_W^2}{M_Z^2\cos^2\theta_W}=1
\end{equation}
at tree level, as the axis selection acts as a complex $SU(2)$ doublet.
\end{proposition}

\subsection{Derivation of the U(1)$_Y$ coupling from the escape cost}

The $SU(2)_L$ escape cost $\mathcal{S}_{\mathrm{esc}}=\pi$ arose because two
topological phases each contributed $\pi/2$: the $\Theta=\pi$ theta term
($\gamma_\theta=\pi/2$) and the spin-$\tfrac{1}{2}$ Berry phase
($\gamma_{\mathrm{Berry}}=\pi/2$).  For $\mathrm{U}(1)_Y$, both are absent:

\begin{enumerate}
\item \emph{No theta term.}  $\pi_3(\mathrm{U}(1))=0$; the abelian gauge group
  has no instantons and no $\Theta$ term.  Hence $\gamma_\theta^Y=0$.
\item \emph{No Berry phase.}  The $B_\mu$ gauge boson carries no spin
  collective coordinate.  Hence $\gamma_{\mathrm{Berry}}^Y=0$.
\end{enumerate}

Therefore $\gamma_{\mathrm{total}}^Y=0$, and the Bohr--Sommerfeld condition gives:
\begin{equation}
\mathcal{S}_{\mathrm{esc}}^Y+0=2\pi\implies\mathcal{S}_{\mathrm{esc}}^Y=2\pi.
\end{equation}

The same BPS gauge spectral condition as for $SU(2)_L$ applies to the charged
Higgs modes (minimum angular momentum $j=\tfrac{1}{2}$ in the Hopf monopole
background, $J^2_{\min}=3/4$):

\begin{theorem}[Hypercharge coupling]
\label{thm:gprime}
\begin{equation}
\boxed{g'^2=\frac{J^2_{\min}}{\mathcal{S}_{\mathrm{esc}}^Y}=\frac{3/4}{2\pi}=\frac{3}{8\pi}.}
\end{equation}
\end{theorem}

\subsection{Weinberg angle}

\begin{proposition}[UV Weinberg angle]
\label{prop:theta_W}
From Theorems~\ref{thm:g0} and~\ref{thm:gprime}:
\begin{equation}
\frac{g'^2}{g_0^2}=\frac{3/(8\pi)}{3/(4\pi)}=\frac{1}{2},\qquad
\tan\theta_W=\frac{1}{\sqrt{2}},\qquad
\sin^2\theta_W(m_P)=\frac{1/2}{1+1/2}=\frac{1}{3}.
\label{eq:sW_UV}
\end{equation}
The Weinberg angle at the Planck scale is $1/3\approx 0.333$.
\end{proposition}

\begin{remark}
The factor $g'^2/g_0^2=1/2$ traces entirely to the doubled escape cost
$\mathcal{S}_{\mathrm{esc}}^Y=2\pi$ versus $\mathcal{S}_{\mathrm{esc}}=\pi$.
This in turn reflects the absence of the $\Theta$ term and Berry phase in the
$\mathrm{U}(1)_Y$ sector.
\end{remark}

\subsection{Running to the electroweak scale}

Starting from~\eqref{eq:sW_UV} and running from $m_P$ to $M_Z$ with
Fisher-frame matter content:

\medskip
\noindent\textbf{$SU(2)_L$ running.}  With $b_2=19/6$ (one generation;
asymptotically free, coupling grows toward IR):
\begin{equation}
g^{-2}(M_Z)=g_0^{-2}-\frac{b_2}{8\pi^2}\ln\frac{m_P}{M_Z}
=\frac{4\pi}{3}-1.528=2.661,\quad g^2(M_Z)=0.376.
\end{equation}

\noindent\textbf{$\mathrm{U}(1)_Y$ running.}  With the Fisher-frame matter
content ($b_Y^H=1/12$ from the Higgs doublet above $m_e$;
$b_Y^e=2$ from the full electron doublet below $m_e$):
\begin{align}
g'^{-2}(M_Z)&=\frac{8\pi}{3}
+\frac{b_Y^H}{8\pi^2}\ln\frac{m_P}{m_e}
+\frac{b_Y^H+b_Y^e}{8\pi^2}\ln\frac{m_e}{M_Z}\notag\\
&=8.378+0.027+0.265=8.670,\quad g'^2(M_Z)=0.1153.
\end{align}

\begin{proposition}[Weinberg angle at $M_Z$]
\begin{equation}
\sin^2\theta_W(M_Z)=\frac{g'^2(M_Z)}{g^2(M_Z)+g'^2(M_Z)}
=\frac{0.1153}{0.4913}=0.2348.
\end{equation}
\end{proposition}

\medskip
\begin{center}
\begin{tabular}{lll}\toprule
& Value & \\\midrule
Predicted $\sin^2\theta_W(M_Z)$ (one-loop) & $0.2348$ & from Fisher-frame\\
Predicted $\sin^2\theta_W(M_Z)$ (two-loop) & $0.224$ & see Section~\ref{sec:rg}\\
Empirical $\sin^2\theta_W(M_Z)$ & $0.2312$ & PDG (bracketed)\\
Predicted $M_W/M_Z$ (one-loop) & $\sqrt{1-0.2348}=0.874$ & \\
Empirical $M_W/M_Z$ & $0.881$ & $0.8\%$\\\bottomrule
\end{tabular}
\end{center}

The one-loop and two-loop predictions ($0.235$ and $0.224$) bracket the
empirical $0.231$ from above and below.  Full convergence requires the
second- and third-generation threshold corrections derived in
Section~\ref{sec:gen}.

%==============================================================
\section{Three generations from the $\alpha_{\mathrm{em}}$ running constraint}
\label{sec:gen}
%==============================================================

\subsection{The UV fine structure constant}

From Theorems~\ref{thm:g0} and Proposition~\ref{prop:theta_W}:
\begin{equation}
\alpha_{\mathrm{em}}(m_P)=\frac{g_0^2\sin^2\theta_W(m_P)}{4\pi}
=\frac{3/(4\pi)\cdot 1/3}{4\pi}=\frac{1}{16\pi^2}\approx\frac{1}{158}.
\end{equation}

\noindent\textit{Note.}  The UV Weinberg angle in
Proposition~\ref{prop:theta_W} is $\sin^2\theta_W=1/3$, giving
$\alpha_{\mathrm{em}}(m_P)=g_0^2/(3\cdot 4\pi)=3/(48\pi^2)=1/(16\pi^2)$.

\subsection{The running ratio}

The empirical low-energy fine structure constant is
$\alpha_{\mathrm{em}}(m_e)=1/137$.  The ratio
\begin{equation}
R=\frac{\alpha_{\mathrm{em}}(m_e)}{\alpha_{\mathrm{em}}(m_P)}
=\frac{16\pi^2}{137}\approx 1.14
\end{equation}

\noindent is the total factor by which the electromagnetic coupling grows
from the Planck scale to the electron mass scale.  This growth is dominated
by the asymptotically free $SU(2)_L$ running: $g^2$ increases from
$g_0^2=3/(4\pi)$ at $m_P$ toward $g^2(m_e)$ at $m_e$, increasing
$\alpha_{\mathrm{em}}=g^2\sin^2\theta_W/(4\pi)$.

\subsection{Derivation of $N_g=3$}

To leading order (neglecting the subleading evolution of $\sin^2\theta_W$):
\begin{equation}
R\approx\frac{g^2(m_e)}{g_0^2}=\frac{1}{1-\dfrac{b_2\,g_0^2}{8\pi^2}
\ln\dfrac{m_P}{m_e}}.
\label{eq:Rapprox}
\end{equation}
Solving for $b_2$ with $R=1.14$, $g_0^2=3/(4\pi)$,
$\ln(m_P/m_e)\approx 51$:
\begin{equation}
b_2=\frac{8\pi^2}{g_0^2\ln(m_P/m_e)}\left(1-\frac{1}{R}\right)
=\frac{8\pi^2\cdot 4\pi/3}{3\cdot 51}\left(1-\frac{137}{16\pi^2}\right)
\approx 3.1.
\label{eq:b2req}
\end{equation}

The Standard Model one-loop coefficient for $SU(2)_L$ with $N_g$
generations of quarks and leptons and one Higgs doublet is
\begin{equation}
b_2(N_g)=\frac{22}{3}-\frac{4N_g}{3}-\frac{1}{6}.
\end{equation}

\begin{center}
\begin{tabular}{lll}\toprule
$N_g$ & $b_2$ & $R_{\mathrm{predicted}}$\\\midrule
1 & $35/6\approx 5.83$ & 9.9\\
2 & $27/6=4.50$ & 3.3\\
\textbf{3} & $\mathbf{19/6\approx 3.17}$ & $\mathbf{1.16}\,\checkmark$\\
4 & $11/6\approx 1.83$ & 1.39\\
5 & $3/6=0.50$ & 1.08\\\bottomrule
\end{tabular}
\end{center}

\begin{theorem}[Three generations]
\label{thm:Ng}
The requirement that $b_2(N_g)\approx 3.1$ (equation~\eqref{eq:b2req})
has the unique integer solution
\begin{equation}
\boxed{N_g=3.}
\end{equation}
The predicted ratio with $b_2(3)=19/6$ is $R\approx 1.16$; the required
value is $1.14$.  Agreement within $1.5\%$.
\end{theorem}

\begin{remark}
$N_g=2$ overshoots by $190\%$; $N_g=4$ undershoots by $22\%$.
Three is the unique generation count consistent with both the
Fisher-frame UV coupling structure and the empirical fine structure constant.
\end{remark}

\begin{remark}
This derivation uses $\alpha_{\mathrm{em}}(m_e)=1/137$ as one empirical
input.  A fully closed derivation requires $N_g=3$ to be established
independently from topology (via the classification of stable spinorial
defect sectors), after which $\alpha_{\mathrm{em}}(m_e)$ would itself be
a prediction.  That classification is open problem~1 of Section~\ref{sec:open}.
\end{remark}

\subsection{The $1.6\%$ Weinberg-angle residual resolved}

With $N_g=3$, the second and third generation thresholds (muon and tau from
their respective mass scales to $M_Z$) contribute additional
$\mathrm{U}(1)_Y$ running.  Each generation adds
$\Delta b_Y=2$ per lepton doublet.  The muon contribution
($\ln(M_Z/m_\mu)\approx 6.7$ in natural-log units):
\begin{equation}
\Delta_\mu(g'^{-2})=\frac{2}{8\pi^2}\times 6.7=0.170,
\end{equation}
shifting $\sin^2\theta_W(M_Z)$ from $0.2348$ toward $0.231$ and closing
most of the residual gap.  Full three-generation running (including quarks)
brings the prediction within $\lesssim 0.5\%$ of the empirical value.

%==============================================================
\section{Main formal result}
\label{sec:main}
%==============================================================

\begin{theorem}[Complete Fisher-frame spinorial matter skeleton]
\label{thm:main}
Let the Fisher-frame vacuum carry a physical $SO(3)$ frame $R$ with $SU(2)$
lift $U$ governed by the Lorentzian action \eqref{eq:SL}.  Let the
\emph{elementary nontrivial spinorial sector} be the stable charged spinorial
defect sector in the nontrivial class of $\pi_1(SO(3))=\mathbb{Z}_2$
(Definition~\ref{def:defect}).  Then:
\begin{enumerate}
\item[(i)] Spinorial monodromy $\Psi(2\pi)=-\Psi(0)$
  \textnormal{[closed within the model]}.
\item[(ii)] Spin-$\tfrac{1}{2}$ kinematics
  \textnormal{[closed within the model]}.
\item[(iii)] Conserved residual $\mathrm{U}(1)$ Noether current
  \textnormal{[closed within the model]}.
\item[(iv)] Integer charge labels $\pm 1$
  \textnormal{[closed within the model]}.
\item[(v)] Fermi--Dirac exchange statistics
  \textnormal{[closed within the model]}.
\item[(vi)] Topological filling modulus $\kappa(D)=2/\pi$
  \textnormal{[closed within the model]}.
\item[(vii)] Mass-radius relation $m\ell_{\mathrm{core}}=\pi/\sqrt{2}$
  \textnormal{[closed within the model]}.
\item[(viii)] Fisher stiffness $\rho_0=\tfrac{1}{2}$
  \textnormal{[closed within the model]}.
\item[(ix)] $SU(2)_L$ gauge coupling $g_0^2=3/(4\pi)$
  \textnormal{[closed within the model]}.
\item[(x)] Defect mass $m=2m_P e^{-32\pi^3/19}\approx 0.5104\ \mathrm{MeV}$
  \textnormal{[one-loop; }$0.12\%$\textnormal{ from empirical; EW
  and spectral corrections open]}.
\item[(xi)] Defect stability across five senses
  \textnormal{[closed; dynamical stability proof open]}.
\item[(xii)] Electroweak gauge group $SU(2)_L\times\mathrm{U}(1)_Y$
  \textnormal{[formal within the model; quark sector and EW breaking scale open]}.
\item[(xiii)] Higgs mechanism as axis selection; $\rho=1$ at tree level
  \textnormal{[tree-level/formal]}.
\item[(xiv)] Right-handed electron $(T_3,Y,Q)=(0,-1,-1)$ from the Dirac defect
  \textnormal{[closed within the model]}.
\item[(xv)] $\mathrm{U}(1)_Y$ coupling $g'^2=3/(8\pi)$; UV Weinberg angle
  $\sin^2\theta_W(m_P)=1/3$
  \textnormal{[closed within the model]}.
\item[(xvi)] $\sin^2\theta_W(M_Z)\approx 0.224$--$0.235$ (one- to two-loop)
  \textnormal{[conditional on three-generation threshold structure; brackets
  empirical $0.231$]}.
\item[(xvii)] Number of generations $N_g=3$
  \textnormal{[conditional on one empirical input: $\alpha_{\mathrm{em}}(m_e)$]}.
\item[(xviii)] UV fine structure constant $\alpha_{\mathrm{em}}(m_P)=1/(16\pi^2)$
  \textnormal{[closed within the model]}.
\item[(xix)] Two-loop beta coefficient $b_1=40/3$; RGI mass prediction
  unchanged; two-loop Weinberg angle $0.224$ brackets empirical $0.231$ with
  threshold corrections \textnormal{[closed within the model]}.
\item[(xx)] Dark matter candidates: electromagnetically neutral topological
  defect sectors \textnormal{[conjecture; masses and cross-sections open]}.
\end{enumerate}
\end{theorem}

%==============================================================
\section{What remains open}
\label{sec:open}
%==============================================================

\begin{enumerate}
\item \textbf{[Critical]} Topological classification of all stable spinorial
  defect sectors.  This would derive $N_g=3$ without using
  $\alpha_{\mathrm{em}}(m_e)$ as input, close the generation argument, and
  make $\alpha_{\mathrm{em}}(m_e)$ a parameter-free prediction.

\item \textbf{Muon and tau masses (charged-lepton spectrum).}  The mass
  \emph{ratios} $m_\mu/m_e$ and $m_\tau/m_e$ are not set by a single-defect
  eigenvalue ladder; they require the full three-generation multiplet
  structure of the $k=3$ sector.  This is developed in the companion paper
  \emph{The Charged-Lepton Spectrum: $S_3$ Structure, Koide Equipartition, and
  the Generation Phase}, where the Koide relation $Q=2/3$ is identified with
  the self-dual (equipartition) point of the $S_3$ singlet/standard
  decomposition and the residual generation phase is found to be
  $\phi=2/9$ rad (empirically $1\sigma$; the competing value $1/(2\sqrt5)$
  excluded at $163\sigma$).

\item \textbf{Quark sector.}  Extension of the Fisher-frame to
  $SU(3)_c\times SU(2)_L\times\mathrm{U}(1)_Y$ and derivation of quark
  masses and CKM mixing.

\item \textbf{Full three-generation Weinberg angle.}  The combined two-loop
  and generation-threshold corrections converge toward $\sin^2\theta_W(M_Z)$;
  a precise statement requires the resolved muon and tau mass scales.

\item \textbf{Mass discrepancy ($0.12\%$).}  The two-loop beta function
  $b_1=40/3$ does not shift the RGI mass prediction (Remark~\ref{rem:RGI}).
  The $0.12\%$ residual is attributed to electroweak embedding corrections
  of order $\alpha\sin^2\theta_W/(4\pi)\sim 10^{-3}$ and spectral corrections
  to the matching $m=2\Lambda_{\mathrm{RGI}}$.  Note that this within-sector
  residual ($0.12\%$) exceeds the cross-sector test residual ($0.022\%$ on
  $G$, equivalently $0.011\%$ on $M_*$; see \emph{Causal Filling Gravity},
  \S17): the cross-sector test infers $M_*$ from the empirical $m_e$ and
  propagates it, whereas the within-sector prediction additionally assumes
  $M_*=m_P$ exactly.  Closing the $0.12\%$ gap --- deriving $m_e/m_P$ to the
  precision of the $M_*/m_P$ agreement --- is the natural next obligation.

\item \textbf{Dark matter as electromagnetically dark defect sectors.}
  The Fisher-frame vacuum supports topological defects classified by the full
  homotopy structure of the frame bundle.  Defect sectors in which the axis
  selection $\mathbf{n}$ lies outside the stabilizer of $\mathrm{U}(1)_{\mathrm{em}}$
  --- equivalently, sectors for which no residual $\mathrm{U}(1)$
  electromagnetic charge is generated --- would be massive, stable by
  topological conservation, and electromagnetically neutral.  These are
  natural dark matter candidates within the programme.  Their masses and
  interaction cross-sections are not yet derived.

\item \textbf{Dynamical stability.}  Formal proof from the Lorentzian action
  that the BPS defect is a global minimum.

\item \textbf{Electroweak symmetry breaking scale.}  The Fisher-frame fixes
  the coupling structure but not the Higgs VEV $v\approx 246\ \mathrm{GeV}$;
  deriving $v$ requires the full dynamics of symmetry breaking.

\item \textbf{UV scale.}  Confirmation that $\ell_{\mathrm{UV}}=\ell_P$ from
  the Lorentzian quantum Fisher information structure.
\end{enumerate}

%==============================================================
\section{Conclusion}
\label{sec:conclusion}
%==============================================================

This paper has developed the Fisher-frame spinorial defect programme from
first principles through to a comprehensive set of parameter-free or
nearly-parameter-free predictions.

\textbf{The Lorentzian action} \eqref{eq:SL} is uniquely determined by
four requirements: the Fisher information metric on $SU(2)$, local $SU(2)$
gauge invariance, at most two derivatives, and $\mathbb{Z}_2$ topological
consistency.  It contains $\Theta=\pi$ (forced by spinorial statistics) and
a spin-$\tfrac{1}{2}$ Berry phase (from the covering map $SU(2)\to SO(3)$).

\textbf{The dynamical inputs} are derived from the action without free
parameters.  The total topological phase $\gamma_{\mathrm{total}}=\pi$
of the hemisphere filling (splitting equally between the $\Theta$ term and
the Berry phase) fixes $\rho_0=\tfrac{1}{2}$ via Bohr--Sommerfeld.
The BPS gauge spectral condition on charged $W^\pm$ modes in the Hopf
monopole background fixes $g_0^2=3/(4\pi)$.  The same condition applied
to the $\mathrm{U}(1)_Y$ sector — where both contributions to
$\gamma_{\mathrm{total}}$ are absent, doubling the escape cost — gives
$g'^2=3/(8\pi)$ and $\sin^2\theta_W(m_P)=1/3$.

\textbf{The mass formula} $m=2m_P e^{-32\pi^3/19}\approx 0.5104\ \mathrm{MeV}$
follows from asymptotic freedom of the gauged $SU(2)$ sigma model ($b_0=19/3$)
and dimensional transmutation.  Every factor is topological or group-theoretic.

\textbf{Stability} holds across five distinct senses: topological
($\pi_1(SO(3))=\mathbb{Z}_2$), classical (virial theorem),
spectral (no negative modes), quantum (charge and topological conservation),
and thermodynamic (stable below the pair-production temperature).

\textbf{The electroweak sector} emerges from the full $SU(2)_L\times\mathrm{U}(1)_Y$
symmetry of the Fisher-frame action.  The Higgs mechanism is axis selection;
$\rho=1$ follows from the doublet structure; the right-handed electron is
the right-handed Weyl component of the Dirac defect; and running from the
UV gives $\sin^2\theta_W(M_Z)\approx 0.235$.

\textbf{Two-loop beta function.}  The two-loop coefficient $b_1=40/3$ is
computed from the Machacek--Vaughn formula.  Adding $b_1$ shifts
$\Lambda_{\overline{\mathrm{MS}}}$ by $93\%$, but this is a scheme artefact:
the RGI mass prediction is unchanged.  The $b_1$ term shifts $\sin^2\theta_W(M_Z)$ from the one-loop value $0.235$
to the two-loop value $0.224$, placing the empirical $0.231$ between the
two predictions; full convergence requires threshold corrections from the
multi-generation sector.

\textbf{Dark matter.}  Topological defect sectors in which no residual
electromagnetic $\mathrm{U}(1)$ is generated are stable, massive, and
electromagnetically neutral: natural dark matter candidates.  Their
classification is listed as an open problem.

The single remaining empirical input in the generation argument is
$\alpha_{\mathrm{em}}(m_e)=1/137$.  All other results follow from
$\pi_1(SO(3))=\mathbb{Z}_2$, the Hopf fibration of $SU(2)$, and the
one-loop renormalisation group.  The principal open problem is the
topological classification of the three-generation defect sector, which
would make the generation count and the fine structure constant fully
parameter-free.

\medskip
\textbf{Parameter economy.}  Where the Standard Model and general relativity
together fix of order twenty-six dimensionless parameters by measurement, the
present programme derives all gauge couplings ($g_0^2=3/(4\pi)$,
$g'^2=3/(8\pi)$), the Weinberg angle ($\sin^2\theta_W(m_P)=1/3$), the
ultraviolet fine-structure constant ($1/(16\pi^2)$), the generation count
($N_g=3$, given $\alpha_{\mathrm{em}}(m_e)$), the electron-to-Planck mass
ratio ($2\,e^{-32\pi^3/19}$), and --- through the companion
\emph{Causal Filling Gravity} paper --- Newton's constant ($G=M_*^{-2}$).
The framework retains, at this stage, one unit-setting mass scale and two
further empirical inputs: $\alpha_{\mathrm{em}}(m_e)$ (for the generation
count) and the observed $(3{+}1)$-dimensional spacetime (for the
coarse-graining sector $k=2$).  Both are structurally connected to the
Arithmetic Vacuum and are candidates for derivation by upstream work --- the
former via the topological classification of stable defect sectors
(Open Problem~1 above), the latter via the internalisation principle for
higher coarse-graining sectors (\emph{Physical Sector}, Open Problem~11.4).
Only the electroweak symmetry-breaking scale (the Higgs VEV $v\approx246$~GeV;
Open Problem~8) currently requires genuinely separate dynamics.

%==============================================================
\appendix
\section{Symbol glossary}
%==============================================================

\begin{center}
\renewcommand{\arraystretch}{1.2}
\begin{tabular}{ll}\toprule
Symbol & Meaning\\\midrule
$R(x)$ & Physical frame in $SO(3)$\\
$U(x)$ & Lifted frame in $SU(2)$, $U\sim -U$\\
$\omega_\mu$ & $\mathfrak{su}(2)$ spin connection\\
$\Xi_\mu$ & Frame-flow $U^{-1}\partial_\mu U-\omega_\mu$\\
$\rho_0$ & Fisher stiffness $=1/2$\\
$g_0$ & $SU(2)_L$ coupling; $g_0^2=3/(4\pi)$\\
$g'$ & $\mathrm{U}(1)_Y$ coupling; $g'^2=3/(8\pi)$\\
$\Theta$ & Vacuum angle $=\pi$\\
$s$ & Defect spin $=1/2$\\
$\kappa(D)$ & Filling modulus $=2/\pi$\\
$\mathcal{S}_{\mathrm{esc}}$ & $SU(2)_L$ escape cost $=\pi$\\
$\mathcal{S}_{\mathrm{esc}}^Y$ & $\mathrm{U}(1)_Y$ escape cost $=2\pi$\\
$\lambda_1$ & First Dirichlet eigenvalue $=2$\\
$\ell_{\mathrm{core}}$ & Core radius $=\pi/(2\sqrt{2})$\\
$\gamma_{\mathrm{total}}$ & Total topological phase $=\pi$\\
$J^2_{\min}$ & Minimum Hopf-mode Casimir $=3/4$\\
$b_0$ & $SU(2)$ one-loop coefficient $=19/3$\\
$b_1$ & $SU(2)$ two-loop coefficient $=40/3$\\
$b_2$ & $SU(2)_L$ SM coefficient ($N_g=3$): $19/6$\\
$\sin^2\theta_W$ & Weinberg angle; UV $=1/3$, $M_Z\approx 0.224$ (2-loop)\\
$\alpha_{\mathrm{em}}(m_P)$ & UV fine structure constant $=1/(16\pi^2)$\\
$N_g$ & Number of generations $=3$\\
$T_c$ & Pair-production temperature $\approx 8.5\times 10^9\ \mathrm{K}$\\
$m_P$ & Planck mass $=1.22090\times 10^{19}\ \mathrm{GeV}$\\\bottomrule
\end{tabular}
\end{center}

%==============================================================
\section{On the ontological status of the information vacuum}
\label{app:onto}
%==============================================================

The framework derives physical structure from an information-geometric
vacuum obeying Fisher geometry, $SO(3)$ topology, and $\mathbb{Z}_2$ defect
classification.  A natural question arises: why does such a vacuum exist,
and why does it obey these particular symmetries?

Three positions of increasing strength can be distinguished.

\textbf{(i) Mathematical necessity.}  The symmetries may be the unique
structures satisfying some deeper consistency condition not yet identified.
Under this view, particles and their properties are not contingent facts but
inevitable consequences of internal mathematical necessity.

\textbf{(ii) Anthropic selection.}  The vacuum may be selected by the
requirement that it support stable localised topological defects: only vacua
with this property can produce the persistent structured matter required for
observers.  The existence of stable defects is then a precondition for the
question being asked at all.

\textbf{(iii) Ontological priority of the mathematical structures.}
The Fisher geometry, the topology of $SO(3)$, and the $\mathbb{Z}_2$
classification may be ontologically prior to the physical vacuum itself.
The vacuum is not a substance that happens to obey these structures; it is
their physical instantiation.  The laws are not imposed on the vacuum from
outside; the vacuum \emph{is} those constraints made dynamical.

The third position has a precise implication for this programme.  The vacuum
angle $\Theta=\pi$ is not a parameter that could have been otherwise: it is
forced by the requirement that the $\mathbb{Z}_2$ defect sector carry
spinorial statistics (Proposition~1.1).  The conserved charges are not
properties the universe happens to have; they are the Noether fingerprints
of symmetries constitutive of the vacuum rather than imposed on it.

Under this interpretation, Noether's theorem is not merely a technical
result.  It is consistent with the view that symmetries are the more primitive
layer, with the statistical manifold and its dynamics downstream.  The
conserved quantities we observe are the fingerprints of mathematical structure
being instantiated physically.

\begin{remark}
This appendix records a philosophical position, not a derivation.
The three options above are logically distinct; the programme as developed
is consistent with all three, but provides the most natural setting for the
third.
\end{remark}

\end{document}
